% Figure 2. Distinct Floquet mechanisms from committed 60-digit flanks.
\begin{figure*}[t]
\centering
\begin{tikzpicture}[x=1cm,y=1cm,line cap=round,line join=round,
  >=Stealth,every node/.style={font=\footnotesize}]
  \node[font=\small,anchor=north west] at (0.05,4.38) {(a)};
  \begin{scope}[shift={(2.72,1.25)},x=0.55cm,y=0.40cm]
    \path[fill=black!7]
      (-4,6)
      plot[domain=-4:4,samples=100] (\x,{0.25*\x*\x+2})
      -- (0,-2) -- cycle;
    \draw[line width=1.05pt] (0,-2) -- (4,6);
    \draw[line width=1.00pt,dashed] (-4,6) -- (0,-2);
    \draw[line width=1.00pt,dash dot]
      plot[domain=-4:4,samples=100] (\x,{0.25*\x*\x+2});
    \draw[line width=0.75pt] (-4.5,-3.0) rectangle (4.5,7.0);
    \foreach \x in {-4,0,4}
      \draw[line width=0.55pt] (\x,-3.0) -- ++(0,0.20)
        node[below=1pt] {$\x$};
    \foreach \y in {-2,2,6}
      \draw[line width=0.55pt] (-4.5,\y) -- ++(0.16,0)
        node[left=1pt] {$\y$};
    \node[font=\small,anchor=north east] at (4.5,-3.35) {$a$};
    \node[font=\small,rotate=90] at (-5.15,2.0) {$b$};
    \node[font=\scriptsize,anchor=west] at (2.70,3.65) {$G_+=0$};
    \node[font=\scriptsize,anchor=east] at (-2.35,3.60) {$G_-=0$};
    \node[font=\scriptsize,anchor=south] at (-2.55,3.95) {$\Delta=0$};
    \node[font=\scriptsize] at (0,0.45) {stable};
    \fill[black] (0,-2) circle (2.0pt);
    \node[font=\scriptsize,anchor=north west] at (0.15,-2.10) {$(0,-2)$};
    \fill[black] (1.7065,1.4130) circle (2.5pt);
    \node[font=\scriptsize,anchor=east] at (1.48,1.00) {$+1$};
    \fill[black] (1.7007,2.723)
      -- ++(0,0.18) -- ++(0.16,-0.29) -- ++(-0.32,0) -- cycle;
    \node[font=\scriptsize,anchor=west] at (2.00,2.95) {HH};
  \end{scope}
  \begin{scope}[shift={(6.15,0.15)}]
    \node[font=\small,anchor=north west] at (-0.20,4.20) {(b)};
    \draw[line width=0.75pt] (0,0) rectangle (3.35,3.35);
    \foreach \x/\lab in {0.000/0.94,1.436/1.00,2.871/1.06}
      \draw[line width=0.55pt] (\x,-0.08) -- (\x,0)
        node[below=1pt,font=\scriptsize] {$\lab$};
    \foreach \y/\lab in {0.000/-0.06,1.675/0,3.350/0.06}
      \draw[line width=0.55pt] (-0.08,\y) -- (0,\y)
        node[left=1pt,font=\scriptsize] {$\lab$};
    \node[below=8pt,font=\scriptsize] at (1.67,0) {$\mathrm{Re}\,\lambda$};
    \node[rotate=90,anchor=south,font=\scriptsize] at (-0.72,1.67)
      {$\mathrm{Im}\,\lambda$};
    \draw[line width=0.95pt]
      plot[domain=-3.5:3.5,samples=60]
        ({(cos(\x)-0.94)/0.14*3.35},{(sin(\x)+0.06)/0.12*3.35});
    \draw[line width=0.8pt,fill=white] (1.436,1.675) circle (2.1pt);
    \node[font=\scriptsize,anchor=south] at (1.436,1.92) {$+1$};
    \draw[line width=0.9pt,fill=white]
      (0.833,1.675) ++(-0.11,-0.11) rectangle ++(0.22,0.22);
    \draw[line width=0.9pt,fill=white]
      (2.053,1.675) ++(-0.11,-0.11) rectangle ++(0.22,0.22);
    \node[font=\scriptsize,anchor=north] at (0.84,1.48) {real pair};
    \fill[black] (1.428,2.386) circle (2.4pt);
    \fill[black] (1.428,0.964) circle (2.4pt);
    \node[font=\scriptsize,anchor=west] at (1.62,2.55) {unit circle};
  \end{scope}
  \begin{scope}[shift={(10.55,0.15)}]
    \node[font=\small,anchor=north west] at (-0.20,4.20) {(c)};
    \draw[line width=0.75pt] (0,0) rectangle (3.35,3.35);
    \foreach \x/\lab in {0.000/0.38,1.675/0.42,3.350/0.46}
      \draw[line width=0.55pt] (\x,-0.08) -- (\x,0)
        node[below=1pt,font=\scriptsize] {$\lab$};
    \foreach \y/\lab in {0.000/0.86,1.675/0.90,3.350/0.94}
      \draw[line width=0.55pt] (-0.08,\y) -- (0,\y)
        node[left=1pt,font=\scriptsize] {$\lab$};
    \node[below=8pt,font=\scriptsize] at (1.67,0) {$\mathrm{Re}\,\lambda$};
    \node[rotate=90,anchor=south,font=\scriptsize] at (-0.72,1.67)
      {$\mathrm{Im}\,\lambda$};
    \draw[line width=0.95pt]
      plot[domain=61:68,samples=40]
        ({(cos(\x)-0.38)/0.08*3.35},{(sin(\x)-0.86)/0.08*3.35});
    \fill[black] (1.204,2.204) circle (2.6pt);
    \node[font=\scriptsize,anchor=east] at (1.02,2.35) {$+$};
    \draw[line width=0.95pt,fill=white] (2.580,1.558) circle (2.4pt);
    \node[font=\scriptsize,anchor=west] at (2.78,1.70) {$-$};
    \draw[line width=0.9pt,fill=white]
      (1.633,1.172) ++(-0.11,-0.11) rectangle ++(0.22,0.22);
    \draw[line width=0.9pt,fill=white]
      (2.303,2.554) ++(-0.11,-0.11) rectangle ++(0.22,0.22);
    \node[font=\scriptsize,anchor=south] at (2.30,2.78) {quartet};
    \node[font=\scriptsize,anchor=north] at (1.63,0.95) {quartet};
  \end{scope}
\end{tikzpicture}
\caption{\label{fig:floquet}Distinct Floquet mechanisms on the continuation
component. (a)~Stability geometry in the physical $(a,b)$ trace plane.
Boundaries $G_+=0$ (solid), $G_-=0$ (dashed), and $\Delta=0$ (dot-dashed)
correspond to $+1$, $-1$, and mode-collision conditions. The filled circle is
the lower $+1$ transition, the triangle is the upper Hamiltonian--Hopf
transition, and the mixed vertex is $(0,-2)$. (b)~Near $+1$, a real reciprocal
pair (open squares) reaches $+1$ and enters the unit circle (filled circles).
(c)~Two unit-circle modes of opposite Krein signature (filled $+$, open $-$)
leave the unit circle as a reciprocal complex quartet (open squares).}
\end{figure*}
